\documentclass{article}
\usepackage[utf8]{inputenc}
\usepackage[margin=1in,left=1.5in,includefoot]{geometry}
\usepackage{booktabs}
\usepackage{float}
% Header & Footer Stuff

\usepackage{fancyhdr}
\usepackage{wasysym}
\usepackage{amsmath}
\usepackage{tikz}
\usepackage{tabularx}
\usepackage{amssymb}
\usepackage{graphicx}
\pagestyle{fancy}
\lhead{MATH 220/199}
\rhead{April 21, 2026}
\fancyfoot{}
\lfoot{David Gallardo}
\fancyfoot[R]{}
%027601193
% The Main Document
\begin{document}

\begin{center}
 \LARGE\bfseries MATH 220/199 Week 14 Day 1
 \line(1,0){430}
\end{center}
\section*{Reminders}
\begin{enumerate}
    \item Office hours are held in Davenport 330 at 3pm on Thursdays.
    \item Do the Merit Quiz on the Merit Canvas page.
\end{enumerate}
\section*{Scaffolding}
The hardest part of u-sub, like many other calculus techniques, is keeping your work and ideas organized. For the following problems, separate your work into the four boxes given: In box 1, identify $u, du$. In box 2, rewrite the integral in terms of $u, du$. In box 2, perform the integration in terms of $u$. In box 4, find your final answer in terms of $x$.
\begin{enumerate}
    \item $\int \sin^2(x)\cos(x) dx$
    \newline
    \newline
    \begin{tikzpicture}[scale = 0.8]
        \draw (0, 0) -- (16, 0);
        \draw (0, 0) -- (0, -16);
        \draw (0, -16) -- (16, -16);
        \draw (16, 0) -- (16, -16);
        \draw (0, -8) -- (16, -8);
        \draw (8, 0) -- (8, -16);
    \end{tikzpicture}
    \newpage
    \item $\int \frac{x}{\sqrt{1-4x^2}} dx$
    \newline
    \newline
    \begin{tikzpicture}[scale = 0.6]
        \draw (0, 0) -- (20, 0);
        \draw (0, 0) -- (0, -16);
        \draw (0, -16) -- (20, -16);
        \draw (20, 0) -- (20, -16);
        \draw (0, -8) -- (20, -8);
        \draw (10, 0) -- (10, -16);
    \end{tikzpicture}
    \vfill
    \item $\int \cot(x) dx$
    \newline
    \newline
    \begin{tikzpicture}[scale = 0.6]
        \draw (0, 0) -- (20, 0);
        \draw (0, 0) -- (0, -16);
        \draw (0, -16) -- (20, -16);
        \draw (20, 0) -- (20, -16);
        \draw (0, -8) -- (20, -8);
        \draw (10, 0) -- (10, -16);
    \end{tikzpicture}
    \vfill
\end{enumerate}

\end{document}
